\documentclass[11pt]{article}
\newcommand{\DocTitle}{Foundations and Customer-Centered Method}
\newcommand{\DocSubtitle}{Evidence-driven design, research, delivery and governance}
\newcommand{\DocID}{DOS-01}
% design-of-sites-style.tex
% Suite preamble for the Web Design Patterns / Design of Sites suite.
%
% RECONSTRUCTION PROVENANCE: the nine uploaded suite documents load a shared
% preamble with `\input{style.tex}`, but no `style.tex` was present in the
% upload batches or anywhere in this repository's history. This file is a
% documented, suite-specific reconstruction inferred from the interface the
% documents actually use: per-document metadata (\DocTitle, \DocSubtitle,
% \DocID defined before this file is input), the front-matter renderer
% \MakeFront, the provenance marker \Provenance, the callout environments
% groundedbox/modernbox/notebox/decisionbox, and the colors
% Ink/Ocean/Mint/Amber/Faint referenced in document TikZ code. It is NOT the
% recovered original style; the original visual design is unknown.
%
% pdfLaTeX only.
\usepackage[T1]{fontenc}
\usepackage[utf8]{inputenc}
\usepackage{lmodern}
\usepackage{microtype}
\usepackage[letterpaper,margin=0.85in,headheight=15pt]{geometry}
\usepackage[table]{xcolor}
\definecolor{Ink}{HTML}{1E2A36}
\definecolor{Ocean}{HTML}{0F4C72}
\definecolor{Mint}{HTML}{0E7E6B}
\definecolor{Amber}{HTML}{B07A1E}
\definecolor{Faint}{HTML}{F2F7F9}
\usepackage{amsmath,amssymb}
\usepackage{array,booktabs,longtable,tabularx,multirow}
\usepackage{enumitem}
\usepackage[most]{tcolorbox}
\usepackage{listings}
\usepackage{tikz}
\usetikzlibrary{arrows.meta,positioning,shapes.geometric,fit,calc}
\usepackage{needspace}
\usepackage{titlesec}
\usepackage{fancyhdr}
\usepackage{lastpage}
\usepackage{xurl}
\usepackage{hyperref}
\hypersetup{colorlinks=true,linkcolor=Ocean,urlcolor=Mint,citecolor=Ocean,pdftitle={\DocTitle},pdfauthor={Web Design Patterns Engineering Suite}}
\setlength{\parindent}{0pt}
\setlength{\parskip}{5pt plus 1pt}
\setlength{\emergencystretch}{2em}
\renewcommand{\arraystretch}{1.18}
\setlist[itemize]{itemsep=2pt,topsep=3pt,leftmargin=1.5em}
\setlist[enumerate]{itemsep=2pt,topsep=3pt,leftmargin=1.7em}
\titleformat{\section}{\Large\bfseries\color{Ocean}}{\thesection}{0.7em}{}[\vspace{-3pt}\color{Mint}\titlerule]
\titleformat{\subsection}{\large\bfseries\color{Ink}}{\thesubsection}{0.6em}{}
\titleformat{\subsubsection}{\normalsize\bfseries\color{Mint}}{\thesubsubsection}{0.55em}{}
\pagestyle{fancy}
\fancyhf{}
\fancyhead[L]{\sffamily\footnotesize\color{Ocean} Design of Sites | \DocID}
\fancyhead[R]{\sffamily\footnotesize\color{Ink} Web Design Patterns Suite}
\fancyfoot[L]{\sffamily\scriptsize Original interpretive guidance; pattern names used as a taxonomy}
\fancyfoot[R]{\sffamily\scriptsize\thepage\ / \pageref{LastPage}}
\renewcommand{\headrulewidth}{0.3pt}
\lstdefinestyle{dositescode}{basicstyle=\ttfamily\footnotesize,breaklines=true,columns=fullflexible,keepspaces=true,showstringspaces=false,frame=single,rulecolor=\color{Mint},backgroundcolor=\color{Faint},aboveskip=8pt,belowskip=8pt}
\lstset{style=dositescode}
% Callouts. groundedbox marks statements verified against the public
% bibliographic record; modernbox marks newly authored modern engineering
% interpretation; notebox is a study/summary callout; decisionbox records
% decisions and trade-offs.
\newtcolorbox{groundedbox}[1][]{enhanced,breakable,colback=Faint,colframe=Mint,boxrule=0.5pt,arc=1.5mm,title={Grounded in the public record},fonttitle=\bfseries\small,#1}
\newtcolorbox{modernbox}[1][]{enhanced,breakable,colback=white,colframe=Ocean,boxrule=0.5pt,arc=1.5mm,title={Modern engineering interpretation},fonttitle=\bfseries\small,#1}
\newtcolorbox{notebox}[1][]{enhanced,breakable,colback=Faint,colframe=Amber,boxrule=0.5pt,arc=1.5mm,#1}
\newtcolorbox{decisionbox}[1][]{enhanced,breakable,colback=white,colframe=Amber,boxrule=0.5pt,arc=1.5mm,title={Decision record},fonttitle=\bfseries\small,#1}
% Front matter renderer: consumes \DocTitle, \DocSubtitle, \DocID.
\newcommand{\MakeFront}{%
  \thispagestyle{empty}%
  {\color{Mint}\rule{\textwidth}{3pt}}\par\vspace{16pt}
  {\sffamily\small\color{Amber}\bfseries \DocID\par}\vspace{10pt}
  {\Huge\bfseries\color{Ocean} \DocTitle\par}\vspace{12pt}
  {\Large\color{Ink} \DocSubtitle\par}\vspace{14pt}
  {\small\color{Ink} Independent educational and engineering companion to \emph{The Design of Sites} (2nd edition). Original explanatory material; the copyrighted pattern text is not reproduced.\par}
  \vspace{10pt}{\color{Mint}\rule{\textwidth}{1pt}}\par\vspace{8pt}
  \tableofcontents\clearpage}
% Provenance marker used at the top of each document body.
\newcommand{\Provenance}{\par\noindent{\footnotesize\sffamily\color{Amber}\bfseries Provenance:}~{\footnotesize\color{Ink}\itshape claims grounded in the public bibliographic record are boxed and cited; everything else is newly authored interpretation.}\par\smallskip}

\begin{document}
\MakeFront
\section{Purpose, scope and editorial boundary}

\Provenance
\begin{groundedbox}
The publisher confirms five foundational chapters before the 13 pattern groups and five research appendixes. The historical book was published in December 2006 and has 1,024 pages. This guide uses its organization as a taxonomy, not its copyrighted explanations. See [B1--B2].
\end{groundedbox}
\begin{modernbox}
A web property is now simultaneously an interaction surface, a distributed software system, an identity and data-processing system, and an operational service. Successful design therefore needs cross-functional evidence from product research, engineering, security, accessibility and SRE.
\end{modernbox}
\textbf{Audience.} Software engineers, product designers, architects, application security engineers, content strategists and delivery teams.
\textbf{Definition of done.} A user journey is delivered only when its task outcomes, safety properties, accessibility, performance and operational recovery paths have been demonstrated.

\subsection{1. Customer-Centered Web Design: More Than a Good Idea}
\begin{groundedbox}Chapter topic and placement verified in the publisher contents [B1]. The interpretation below is newly authored.\end{groundedbox}
\textbf{Engineering interpretation.} Select a high-value customer task, identify its context, constraints and measurable outcome. Differentiate observed behavior from stakeholder assumptions.

\textbf{Field exercise.} Interview at least five representative participants across relevant abilities, devices and roles; triangulate interview claims with analytics and task observation.

\textbf{Constraint.} Risk: optimizing conversions through coercion or dark patterns. Add explicit trust and consent constraints.

\textbf{Evidence to retain:} problem statement; design alternatives; participants or workload; observation notes; adverse outcomes; decision record; owner and review date.

\subsection{2. Making the Most of Web Design Patterns}
\begin{groundedbox}Chapter topic and placement verified in the publisher contents [B1]. The interpretation below is newly authored.\end{groundedbox}
\textbf{Engineering interpretation.} Capture a recurring problem, design context, trade-offs, candidate response, related patterns and falsifiable success criterion.

\textbf{Field exercise.} Create a pattern registry tied to design-system components, architecture decisions and test evidence.

\textbf{Constraint.} Risk: pattern overgeneralization. Every reuse must evaluate current user and deployment context.

\textbf{Evidence to retain:} problem statement; design alternatives; participants or workload; observation notes; adverse outcomes; decision record; owner and review date.

\subsection{3. Knowing Your Customers: Principles and Techniques}
\begin{groundedbox}Chapter topic and placement verified in the publisher contents [B1]. The interpretation below is newly authored.\end{groundedbox}
\textbf{Engineering interpretation.} Use evidence from contextual inquiry, interviews, usability sessions, field data, accessibility studies and support tickets.

\textbf{Field exercise.} Maintain segment assumptions, recruiting protocol, research consent and a research evidence ledger.

\textbf{Constraint.} Risk: collection and retention of excess participant data. Minimize identifiable research records.

\textbf{Evidence to retain:} problem statement; design alternatives; participants or workload; observation notes; adverse outcomes; decision record; owner and review date.

\subsection{4. Involving Customers with Iterative Design}
\begin{groundedbox}Chapter topic and placement verified in the publisher contents [B1]. The interpretation below is newly authored.\end{groundedbox}
\textbf{Engineering interpretation.} Iterate from low-fidelity sketches to interactive prototypes to production experiments with an explicit hypothesis.

\textbf{Field exercise.} Run moderated task studies and technical spikes before scaling fidelity; record failed hypotheses.

\textbf{Constraint.} Risk: shipping experiments without guardrails. Predefine rollback, privacy and availability limits.

\textbf{Evidence to retain:} problem statement; design alternatives; participants or workload; observation notes; adverse outcomes; decision record; owner and review date.

\subsection{5. Processes for Developing Customer-Centered Sites}
\begin{groundedbox}Chapter topic and placement verified in the publisher contents [B1]. The interpretation below is newly authored.\end{groundedbox}
\textbf{Engineering interpretation.} Treat discovery, exploration, refinement, production, implementation, launch and maintenance as evidence-producing activities, not rigid phase gates.

\textbf{Field exercise.} Map delivery artifacts to backlog stories, automated CI policies, security gates and post-launch monitoring.

\textbf{Constraint.} Risk: waterfall stage-gate misuse. Use iterative cycles with documented feedback loops.

\textbf{Evidence to retain:} problem statement; design alternatives; participants or workload; observation notes; adverse outcomes; decision record; owner and review date.

\section{A continuous delivery loop}

\begin{enumerate}
\item \textbf{Observe:} select critical tasks and baseline abandonment, confusion, errors, latency and accessibility defects.
\item \textbf{Frame:} document user outcomes, jobs-to-be-done, failure cases and security/privacy invariants.
\item \textbf{Design:} select applicable patterns from the registry; prototype and record alternatives.
\item \textbf{Validate:} conduct task testing and boundary testing with representative users and devices.
\item \textbf{Implement:} ship through component contracts, service/API contracts and CI gates.
\item \textbf{Operate:} collect field measurements, support signals and actionable telemetry with data minimization.
\item \textbf{Revise:} review evidence and update the pattern choice or architecture decision.
\end{enumerate}
\begin{decisionbox}
Do not equate more automation with better experience. A fully passing build can still mislead, exclude or obstruct users. Pair machine-checkable quality gates with ongoing observation of real tasks.
\end{decisionbox}

\section{Outcome and constraint scorecard}

\begin{tabularx}{\linewidth}{@{}p{32mm}X X@{}}\toprule
Dimension & Observable measure & Minimum verification practice \\\midrule
Task success & Success per journey, retries, abandonment & Moderated test and funnel data \\
Comprehension & First-click correctness, time to first useful result & Think-aloud and unmoderated studies \\
Accessibility & WCAG 2.2 AA conformance claim by scope & Manual assistive-tech review and automated checks \\
Security & Exploitable authorization/authentication/data flaws & ASVS 5.0 requirements and threat-led tests \\
Performance & LCP, INP, CLS by device at p75 & Real-user monitoring plus lab budgets \\
Operability & Errors, uptime, latency, rollback & Instrumented SLOs, incident drills \\
Privacy & Collection purpose, deletion time, consent validity & Data inventory and removal test \\
\bottomrule\end{tabularx}
\textit{Targets are project-specific except where a named normative or published threshold is cited; see [S1], [S4], [S10].}

\section{Research and iteration protocol}

\begin{enumerate}
\item Write neutral tasks that describe outcomes, not interface controls. Example: ``Find an item that can arrive within three days and estimate the full price.''
\item Sample across experience levels, devices, assistive technologies, languages and risk profiles relevant to the service.
\item Obtain informed consent for observation and any recording. Avoid production secrets and personal financial information in test data.
\item Log task completion, time, recovery behavior and context; label researcher interpretations separately from observations.
\item Categorize problems by severity and reproducibility; associate each finding with a pattern ID and owner.
\item Re-test the correction, include a negative case, and document any residual compromise.
\end{enumerate}
\begin{modernbox}
For ongoing production measurement, combine task outcome events, web-vitals field measurements and distributed traces. Use sampling and pseudonymization as needed. Telemetry must never capture passwords, authorization headers, payment details or free-form sensitive input.
\end{modernbox}

\section{Applicability and anti-patterns}
\begin{itemize}
\item Choose a \emph{modular monolith} and a single frontend when one team can deliver cohesive user journeys without independent deployments.
\item Consider separate frontends only when business capability ownership, release independence and operational maturity justify additional complexity.
\item Avoid blind replication of legacy patterns such as table-based page layout, imperative pop-ups and counting the total number of assets as the primary performance metric.
\item Do not use personalization as a license for unnecessary tracking, inaccessible variation or unreviewable automated choices.
\item Keep a single user-facing experience even when multiple teams and services deliver its parts.
\end{itemize}


\clearpage
% design-of-sites-references.tex
% Shared references section for the Web Design Patterns / Design of Sites
% suite.
%
% RECONSTRUCTION PROVENANCE: the nine uploaded suite documents end with
% `\input{references.tex}`, but no `references.tex` was present in the upload
% batches or anywhere in this repository's history. The documents cite only
% keys [B1] and [B2] (bracketed tokens such as "[B4, F1, H10, ...]" are
% pattern-ID template placeholders, not citations). The two entries below are
% reconstructed from the bibliographic facts the documents themselves state:
% publisher record and bibliographic-database record of the second edition.
% This file is NOT the recovered original reference list.
\clearpage
\section*{References}
\begin{description}[leftmargin=3.2em,style=nextline]
  \item[{[B1]}] Publisher bibliographic record and public table of contents:
    D.~K.~van Duyne, J.~A.~Landay, J.~I.~Hong, \emph{The Design of Sites:
    Patterns, Principles, and Processes for Crafting a Customer-Centered Web
    Experience}, 2nd edition, Prentice Hall, December 2006, 1{,}024 pages.
    Confirms the five foundational chapters, the thirteen pattern groups
    (A--M), the 107 pattern labels and the research appendixes used as this
    suite's taxonomy.
  \item[{[B2]}] Bibliographic-database record (O'Reilly learning-platform
    catalog entry) for the same second edition; corroborates edition,
    authors, publication date and chapter/group structure. The copyrighted
    pattern explanations themselves were not retrieved or reproduced for
    this suite.
\end{description}

\end{document}
